\documentclass[11pt]{article}
\usepackage{cocina}
\clavetrue
% la clave se consulta, no se lee de corrido: compacta
\newcommand{\cara}[1]{\par\vspace{1.2mm}{\rotulo\footnotesize\color{oro}#1}\par\vspace{.6mm}}
\newcommand{\nj}[1]{\textbf{#1.}\ }
\newcommand{\nota}[1]{{\itshape\color{tinta!75} #1}}
\newcommand{\sep}{\quad·\quad}
\newcommand{\f}[2]{\tfrac{#1}{#2}}

\begin{document}\small\setlength{\parskip}{1mm}

% ═════════════════════════════ FORMATO DE 2.º ═════════════════════════════
\begin{ficha}{Clave}{Clave · Ficha 1 · Servir}{1}
\cara{DELANTE · CORTAR REPARTE, GRAPAR JUNTA}
\nj{1} a) $\f34 = \f68$\sep b) $\f13 = \f26$\sep c) $\f14 = \f3{12}$. Hay más trozos, más
pequeños, y la misma pizza. \nota{La pregunta de «¿te llevas más?» es la de la
clase: que alguien diga que sí y se discuta con la regla.}

\nj{2} a) $\f46 = \f23$\sep b) $\f68 = \f34$\sep c) $\f8{12} = \f23$.

\nj{3} $\f12$, $\f24$ y $\f48$ caen en la misma raya: la cuarta de ocho.

\nj{4} a) $\f4{10}$, $\f6{15}$\sep b) $\f{10}{12}$, $\f{15}{18}$\sep c) $\f34$, $\f68$\sep
d) $\f23$, $\f69$\sep e) $\f34$, $\f9{12}$\sep f) $\f25$, $\f6{15}$. \nota{En c), d) y f)
hay que grapar primero y luego cortar.}

\textbf{Reto.} Infinitos: $\f{n}{2n}$. El más corto es $\f12$: el 1 de arriba no se
puede grapar con nada (1 y 2 no tienen divisor común mayor que 1).

\cara{REVERSO · MÁS PEDIDOS}
\nj{1} a) cortar cada cuarto en 2 ($\cdot 2$): $\f58$\sep b) grapar de 4 en 4 ($:4$):
$\f48 = \f12$\sep c) grapar de 2 en 2: $\f46 = \f23$\sep d) cortar cada quinto en 2:
$\f7{10}$\sep e) grapar de 3 en 3: $\f9{12} = \f34$. \nota{En a) y d) el pedido no
existe en la caja de partida: cinco octavos son dos cuartos y medio.}

\nj{2} $\f54$ = una pizza entera y $\f14$; en la regla, entre el 1 y el 2.

\nj{3} a) $\f12 + \f14 = \f24 + \f14 = \f34$\sep b) $\f13 + \f16 = \f26 + \f16 = \f36 = \f12$.

\nj{4} a) $\f34$\sep b) $\f34$\sep c) $\f34$\sep d) $\f23$\sep e) $\f23$\sep f) $\f13$.
En c), de 6 en 6 de una vez: el \textbf{mcd} del 18 y el 24.

\textbf{Reto.} Sí: grapar $\f46$ de 2 en 2 da $\f23$, y cortar cada tercio en 3 da
$\f69$. \nota{Las dos son «$\f23$ disfrazado»: la irreducible es el nombre común.}
\end{ficha}

\vspace{3mm}
\begin{ficha}{Clave}{Clave · Ficha 2 · Comparar}{2}
\cara{DELANTE · ¿QUIÉN SE LLEVA MÁS?}
\nj{1} a) $\f38 < \f58$: trozos iguales, se cuentan\sep b) $\f13 > \f14$: un trozo cada
uno, el tercio es más grande\sep c) $\f34 = \f68 > \f58$: se corta \textbf{una} caja\sep
d) $\f24 = \f6{12}$ y $\f36 = \f6{12}$: iguales.

\nj{2} $\f23 = \f{10}{15}$, $\f35 = \f9{15}$: $\f23 > \f35$, por $\f1{15}$.

\nj{3} a) en 18: $\f{15}{18} > \f{14}{18}$\sep b) en 20: $\f6{20} > \f5{20}$\sep
c) en 36: $\f{16}{36} > \f{15}{36}$\sep d) en 24: $\f{14}{24} < \f{15}{24}$.

\textbf{Reto.} Sí: en cuarentavos, $\f{30}{40} < \f{31}{40} < \f{32}{40}$. Entre
$\f12$ y $\f24$, ninguna: son la misma. \nota{Entre dos fracciones distintas
siempre cabe otra: basta cortar más fino.}

\cara{REVERSO · MÁS MESAS}
\nj{1} a) $\f38 < \f12 < \f58 < \f34$\sep b) $\f13 < \f12 < \f23 < \f56$.

\nj{2} a) B, $>$\sep b) A, $<$\sep c) C, $<$ ($\f4{10}$)\sep d) D, $<$ ($\f9{12}$ y
$\f{10}{12}$)\sep e) E, $=$ (las dos son $\f23$)\sep f) B, $<$.

\nj{3} a) mesa 7, $\f{25}{40}$ contra $\f{24}{40}$: por $\f1{40}$\sep b) Ana,
$\f{16}{24}$ contra $\f{15}{24}$\sep c) la primera, $\f{15}{20}$ contra $\f{14}{20}$.

\textbf{Reto.} Falla: $\f23$ y $\f34$ tienen la misma distancia y no son iguales; y
con $\f12$ y $\f35$ diría que gana $\f12$.
\end{ficha}

\newpage
\begin{ficha}{Clave}{Clave · Ficha 3 · Por teléfono}{3}
\cara{DELANTE · SOBRE EL PAPEL}
\nj{1} a) 56; no llega (el cortador llega a 15)\sep b) $\f78 = \f{49}{56}$,
$\f67 = \f{48}{56}$\sep c) Nicoleta, por $\f1{56}$.

\nj{2} De la \textbf{segunda}; el 48 es $8\cdot 6$, el de abajo de la primera por el de
arriba de la segunda. Los de abajo no hace falta mirarlos: salen iguales ($8\cdot
7 = 7\cdot 8$).

\nj{3} a) $45 > 28$: $\f57 > \f49$\sep b) $15 < 16$: $\f38 < \f25$\sep c) $35 < 36$:
$\f59 < \f47$\sep d) $49 < 50$: $\f7{10} < \f57$.

\textbf{Reto.} Si $a\cdot d = b\cdot c$, en la secuencia común los de arriba salen
iguales: la misma cantidad. $6\cdot 12 = 72 = 9\cdot 8$.

\cara{REVERSO · MÁS ENCARGOS DE LA PASTELERÍA}
\nj{1} A $\f78$ le falta $\f18$ y a $\f67$, $\f17$. Más pequeño es $\f18$: se lleva
más Nicoleta.

\nj{2} a) $72 = 72$: sí\sep b) $60 = 60$: sí\sep c) $36 \neq 35$: no.

\nj{3} a) $\f{10}{15}$\sep b) $\f68$\sep c) $\f{12}{20}$.

\nj{4} $\f23 < \f7{10} < \f57$ ($14 < 15$, $20 < 21$, $50 > 49$).

\nj{5} a) la de Nicoleta ($16 > 15$)\sep b) el escaparate ($50 > 49$)\sep c) Nick ($21 < 22$).

\textbf{Reto.} En la secuencia de 56: $\f57 = \f{40}{56}$ y $\f34 = \f{42}{56}$, así
que $\f{41}{56}$. Con la cruz: $41\cdot 7 = 287 > 280 = 56\cdot 5$ y
$41\cdot 4 = 164 < 168 = 56\cdot 3$.
\end{ficha}

% ═════════════════════════════ FORMATO DE 1.º ═════════════════════════════
\vspace{3mm}
\begin{ficha}{Clave}{Clave · Formato de 1.º}{0}
\cara{FICHA 1 · SERVIR}
\textbf{Delante.} 2a) $\f12 = \f24$\sep 2b) $\f23 = \f46$\sep 3a) $\f24 = \f12$\sep
3b) $\f46 = \f23$. Reto: $\f48$ y $\f6{12}$.
\textbf{Reverso.} 1) a) $\f26$\sep b) $\f6{10}$\sep c) $\f48$\sep d) $\f34$\sep e) $\f12$\sep
f) $\f13$. 2) Las tres caen en el mismo sitio: son la misma pizza. 3) a) $\f48$, sí
es media\sep b) 4 trozos ($\f46$).

\cara{FICHA 2 · COMPARAR}
\textbf{Delante.} 2a) $\f26 < \f56$\sep 2b) $\f12 > \f14$\sep 2c) $\f34 = \f68 > \f58$.
Reto: iguales ($\f24 = \f36 = \f12$).
\textbf{Reverso.} 1) a) $\f12 > \f38$\sep b) $\f34 > \f58$. 2) a) $<$\sep b) $\f46 <
\f56$\sep c) $>$\sep d) $\f9{12} > \f7{12}$. 3) a) iguales\sep b) Ana ($\f68 > \f58$).

\cara{FICHA 3 · POR TELÉFONO}
\textbf{Delante.} 2a) $6 > 5$: $\f25 > \f13$\sep 2b) $15 > 14$: $\f37 > \f25$. Reto: sí,
$2\cdot 9 = 18 = 3\cdot 6$.
\textbf{Reverso.} 1) a) $20 < 21$: $<$\sep b) $25 > 24$: $>$\sep c) $35 < 36$: $<$\sep
d) $25 > 24$: $>$. 2) a) sí\sep b) sí\sep c) no. 3) Nicoleta ($10 > 9$).
\end{ficha}
\end{document}
